\documentclass[11pt,a4paper]{article}
\usepackage[margin=25mm]{geometry}
\usepackage{amsmath,amssymb,amsthm,hyperref}
\setlength{\emergencystretch}{2em}
\newtheorem{lemma}{Lemma}
\newtheorem{theorem}[lemma]{Theorem}
\newtheorem{corollary}[lemma]{Corollary}
\newcommand{\gen}[1]{\langle#1\rangle}
\title{Ambient transport and internal pair deficits}
\author{Human Compression project: L044 continuation}
\date{}
\begin{document}\maketitle
Let $G=\gen{\rho_0,\ldots,\rho_{r-1}}$ be a transitive finite
string action on $\Omega$. Generators are involutions or identities.
Fix $0\le k<r$ with the ambient omission $G_k$ intransitive, and put
\[
 P=G_{<k},\qquad K=G_{>k},\qquad H=G_{\ge k}.
\]
Thus $P$ centralizes $K$, and $G_k=PK$. Call an actual $H$-orbit
\emph{good} when $K$ is transitive on it, and \emph{bad} otherwise.
These names are confined to this note; they express no source status.

\begin{lemma}[A whole-orbit transport]\label{move}
For every good $H$-orbit $O$ there are $p\in P$ and a bad
$H$-orbit $O'$ such that $O^p$ is a $K$-orbit contained in $O'$.
The bijection $O\longrightarrow O^p$ commutes with every marked
generator $\rho_j$, $j>k$.
\end{lemma}
\begin{proof}
The set $O$ is one $K$-orbit. Since $P$ centralizes $K$, its
translates $O^p$ are $K$-orbits, and their union $B$ is the
$G_k=PK$-orbit containing $O$. It is a proper subset of $\Omega$.
If every $O^p$ were itself an $H$-orbit, $\rho_k$ would preserve
each one and hence preserve $B$. Then $B$ would be invariant under
both $G_k$ and $\rho_k$, contrary to transitivity of $G$.
Choose a translate $W=O^p$ which is not an $H$-orbit. Its containing
$H$-orbit $O'$ strictly contains $W$, so $K$ is intransitive on $O'$.
Commutation with all $\rho_j$, $j>k$, follows from $[P,K]=1$.
\end{proof}

An edge of colour $j$ is \emph{ambient eligible} if its endpoints
lie in different $G_j$-orbits, where $G_j$ is the full ambient
omission. Suppose a set of tail colours
$J\subseteq\{k,\ldots,r-1\}$ is assigned to actual $H$-orbits,
with each colour assigned exactly once. For every assigned colour
choose two ambient-eligible edges in its assigned orbit.

\begin{theorem}[Simultaneous assignment normalization]\label{normal}
The assigned pairs can be replaced by pairs of the same colours
so that every colour is assigned to a bad $H$-orbit. Each new pair
remains ambient eligible. In particular it is eligible for the
full local omission on its new orbit.
\end{theorem}
\begin{proof}
There is no eligible $k$-edge in a good orbit: $K\le G_k$ is
already transitive there. Thus a colour $j$ assigned to a good
orbit satisfies $j>k$. Choose the element $p$ from Lemma~\ref{move}
for that orbit and transport both selected edges by $p$.
They remain distinct actual $j$-edges, since $p$ commutes with
$\rho_j$. Also $P\le G_j$, so $p$ preserves each $G_j$-orbit;
the transported endpoints remain in different such orbits.
Both edges lie in the same bad target orbit $O'$.

Do this for all colours assigned to good orbits, leaving every
original bad-orbit assignment in place. A colour still has just
one assignment, and no assigned pair is subsequently moved away
from a bad target. Different colours cannot cause a conflict in
this construction. Finally, the local omission on $O'$ is the
restriction of $\gen{\rho_i:i\ge k,\ i\ne j}\le G_j$.
Distinct ambient omission orbits therefore imply distinct local
omission orbits. This uses a subgroup inclusion, not equality of
the two orbit partitions.
\end{proof}

\begin{corollary}[The first-colour deficit]\label{deficit}
Suppose all tail indices are assigned, so $J=\{k,\ldots,r-1\}$.
Use the normalized assignments, and let $F_a$ be the colours
assigned to the nontrivial $H$-orbit $O_a$, of size $m_a$.
Write $d_a=|F_a|-m_a/2$, and let $p$ be the unique index with
$k\in F_p$. Then
\[
 d_p\le0,\qquad
 d_a\le
 \begin{cases}-1,&m_a\text{ even},\\-\tfrac12,&m_a\text{ odd}
 \end{cases}\quad(a\ne p).                         \tag{1}
\]
An empty assignment has its exact deficit $d_a=-m_a/2$.
\end{corollary}
\begin{proof}
All full local two-edge hypotheses follow from Theorem~\ref{normal}.
The accepted selected-colour theorem gives $d_p\le0$.
For $a\ne p$ with $F_a\ne\varnothing$, the orbit is bad and
the selected subgroup lies in its intransitive full local omission
of $\rho_k$. Hence $2|F_a|\le m_a-1$. Integrality gives (1).
If $F_a$ is empty, its exact deficit implies the same bounds because
$m_a\ge2$. No transitivity of a selected or projected subgroup is
being assumed.
\end{proof}

\begin{theorem}[An internal centralizer and its exact boundary]\label{internal}
Let $O$ be an actual $H$-orbit of size $2q$, where $q\ge3$.
Suppose an involution $\tau$ centralizes $H$, preserves $O$ and
acts nontrivially on it. For a set $F\subseteq\{k,\ldots,r-1\}$,
suppose $\tau\in G_j$ and at least one ambient-eligible $j$-edge
lies in $O$ for every $j\in F$. Then
\[
                        |F|\le q-1.                  \tag{2}
\]
In equality, the pair images of the colours in $F$ are exactly
the transpositions of a path on the $q$ pairs of $\tau$ in $O$.
Their indices are consecutive; every other tail generator has
trivial pair image.

Suppose additionally that $k\notin F$ and $K$ is intransitive on
$O$. In equality $K|_O$ is the diagonal $S_q$ on two natural rows,
and every tail mark outside $F\cup\{k\}$ fixes $O$ pointwise.
After row labelling, $\rho_k|_O$ flips either all pairs, just the
first pair, or all pairs except the first. The last two alternatives
require that the first path label is $k+1$.
\end{theorem}
\begin{proof}
An element centralizing a transitive group which fixes one point
fixes every point. Thus $\tau|_O$ is fixed-point-free and has
$q$ pairs, preserved by $H$. Choose one ambient-eligible edge for
each $j\in F$ and project it to the pair set. Its endpoints are
distinct pairs because $\tau\in G_j$. The projected edges form a
forest. Indeed an alternate projected path avoiding colour $j$
lifts to a path using the other selected generators and $\tau$,
all in $G_j$, contradicting eligibility of the chosen $j$-edge.
This also excludes repeated projected edges with distinct colours.
The forest bound proves (2).

Suppose equality holds. The forest is a spanning tree $T$.
For each $j\in F$, removing its edge gives two subtrees.
The selected edges of each subtree, together with $\tau$, connect
all points above that subtree inside one $G_j$-orbit. The two
subtrees lie in different $G_j$-orbits, by the chosen edge's
eligibility. Consequently these are exactly the intersections
with $O$ of the two $G_j$-orbits which meet $O$.

Every other tail generator preserves $O$, its pairs, and both
of these cut parts. A colour $i\in F$ preserves every tree cut
except possibly its own, so its pair image fixes all vertices
other than the endpoints of the $i$-edge. It exchanges those
endpoints, hence is exactly that transposition. A tail colour
outside $F$ preserves every cut and fixes all pair vertices.
The image of $H$ on pairs is therefore the full $S_q$.

Incident edges of $T$ give noncommuting transpositions, so their
indices must differ by one. Three edges cannot meet one vertex,
because three distinct integers cannot be pairwise consecutive.
Thus the tree is a path. Its successive indices differ by one
and cannot reverse direction without repeating an index.
The indices are consecutive in one of the two path orientations.

Now suppose the additional hypotheses hold. Since $K$ contains
all marks in $F$, it is transitive on the pair set. Its assumed
intransitivity on points gives exactly two orbits of size $q$,
each meeting each pair once. A pair-trivial element of $K$ fixes
each row and hence every point. Thus $K|_O$ is faithfully the
diagonal $S_q$, and the marks outside $F\cup\{k\}$ are identities
there. Relabel the two rows accordingly.

The mark $\rho_k$ has trivial pair image, so is a vector of pair
flips. If $k+1\notin F$, it commutes with all path transpositions
and the vector is constant. If $k+1\in F$, the path begins there;
commutation with the remaining path fixes a common value on the
other $q-1$ pairs, with an independent value on its first pair.
The zero vector would leave $H$ intransitive on $O$. The nonzero
vectors are exactly those stated.
\end{proof}

For odd $q$, the three displayed flip vectors have respectively
odd, odd and even parity on $O$. In particular the last local
configuration cannot be discarded by parity alone. The theorem
proves deficit $|F|-|O|/2\le-1$ and describes equality; it does
not assert the stronger deficit $-2$. The internal source case
has odd $\tau|_O$, which gives $q$ odd. A two-point orbit has no
ambient-eligible edge of a colour whose omission contains $\tau$,
so its assigned set is empty and its deficit is exactly $-1$.

\paragraph{Source and exact boundary.}
The orbit-transfer strategy is borrowed from
Cameron--Fernandes--Leemans--Mixer, \emph{Highest rank of a polytope
for $A_n$}, arXiv:1605.09173v1, Proposition 5.27, physical page 36,
and the first-colour deficit is used in Proposition 5.31 on page 42.
The proof here specifies the transported whole $K$-orbit and the
proper ambient $G_k$-orbit which forces a bad target. The internal
pair forest and its equality proof supply a separate interface
for the internally acting central involution in Proposition 5.31,
cases (2)--(3), physical pages 42--43. They use only actual ambient
omission cuts and the elementary tree argument, not a projected IP
or an inferred classification of an arbitrary symmetric marking.
The entire accepted selected-colour proof and its dependencies remain
paid. No intersection property, parity or finite computation is used.

Crucially, the chosen pairs in Theorem~\ref{normal} are
\emph{ambient} eligible. Two edges eligible only for a local tail
omission have not been shown to remain eligible after transport
by $P$. Nor is a pair split between two different $H$-orbits an
input pair for this theorem. Applying the note requires the stated
assignment to be supplied; it is not inferred from a bare local
two-fracture condition. This does not establish the stronger internal
odd-orbit deficits or their equality classification in source
Proposition 5.31(2)--(3), and does not complete L044.
\end{document}
